% !TEX root = ../graph-theory.tex
% =====================================================================
%  组合优化 · st-编号
% =====================================================================
\lecture{st-编号：定义与前缀、后缀的连通性，DFS 树的树边与返回边，dfn 与 low 值，求 low 值的算法，“老点、新点”的 st-编号算法及其完整证明，割点}
\label{lec:st-numbering}

设 $G$ 是 $2$-连通图，$st$ 是一条边。$G$ 的 \term{st-编号}是把顶点编为 $1,\dots,n$，使 $s$ 编为 $1$、$t$ 编为 $n$，并且其余每个顶点都有一个编号比它小的邻点和一个编号比它大的邻点。直观上，它把图“拉直”成从 $s$ 到 $t$ 的一列，每个顶点往前、往后都连得出去。平面性检验等算法要用到它。这一讲的算法建立在一次 DFS 的基础上，先讲 DFS 树的几个性质和 low 值，再讲算法和它的证明。

\begin{theorem}\label{thm:st-numbering:prefix}
（讲义定理 2.13）若 $v_1,\dots,v_n$ 是 st-编号（$v_i$ 编号为 $i$），则对每个 $1\leqslant i<n$，$G[\{v_1,\dots,v_i\}]$ 与 $G[\{v_{i+1},\dots,v_n\}]$ 都连通。
\end{theorem}
\begin{proof}
对 $j\leqslant i$，从 $v_j$ 出发，不断走到一个编号更小的邻点（$j\geqslant2$ 时存在），编号严格下降，最终到达 $v_1$，整条路都在前 $i$ 个顶点中。所以前缀中每个顶点都与 $v_1$ 连通。后缀同理，走编号更大的邻点到达 $v_n$。
\end{proof}

\topic{一、DFS 树：树边与返回边}

本讲的 DFS 从边 $ts$ 开始：$t$ 是根，$s$ 是 $t$ 的孩子。记 \texttt{dfn[v]} 为 $v$ 被访问（进栈）的次序，$\texttt{dfn}[t]=0$，$\texttt{dfn}[s]=1$，\texttt{parent[v]} 为 $v$ 在 DFS 树 $T$ 中的双亲，$T_v$ 为以 $v$ 为根的子树。$T$ 中的边叫\term{树边}，$G$ 中其余的边叫\term{返回边}。

\begin{proposition}\label{prop:st-numbering:dfs-tree}
（讲义 2.8 节的性质 (a)--(d)，改正后）设 $G$ 是 $2$-连通图。
\begin{enumerate}
\item[(a)] $t$ 在 $T$ 中只有一个孩子 $s$。
\item[(b)] 若 $u$ 是 $v$ 的祖先，则 $\texttt{dfn}[u]<\texttt{dfn}[v]$；$T_v$ 的顶点的 \texttt{dfn} 是连续的一段，从 $\texttt{dfn}[v]$ 开始。
\item[(c)] 每条返回边连接一个顶点与它的一个祖先。特别地，$u$ 的两个孩子 $v,w$ 的子树 $T_v,T_w$ 之间没有边。
\item[(d)] 若 $u\neq t,s$，$v$ 是 $u$ 的孩子，则有一条返回边从 $T_v$ 连到 $u$ 的某个真祖先。若 $v$ 是 $s$ 的孩子，则有返回边从 $T_v$ 连到 $t$。
\end{enumerate}
\end{proposition}
\begin{proof}
无向图的 DFS 可以看成把每条边换成两条反向的弧之后的有向 DFS，所以上一讲的白路引理在这里同样成立。

(a)：$t$ 进入 $s$ 时，除 $t$ 外的顶点都没有被访问。$G-t$ 连通（$t$ 不是割点），所以从 $s$ 到其余每个顶点都有一条不经过 $t$ 的路，路上的顶点都还没被访问。由白路引理，它们都在 $s$ 出栈之前被访问，都在 $T_s$ 中；$t$ 没有别的孩子。

(b)：祖先先被访问。$v$ 进栈后、出栈前进栈的顶点恰是 $T_v$ 的顶点（白路引理的后一句），它们的 \texttt{dfn} 连续地接在 $\texttt{dfn}[v]$ 之后。

(c)：设 $xy$ 是返回边，不妨设 $x$ 先被访问。$x$ 进栈时 $y$ 还没被访问，边 $xy$ 本身就是一条“白路”，由白路引理 $y$ 在 $x$ 出栈前被访问，所以 $y\in T_x$，$x$ 是 $y$ 的祖先。$T_v$ 与 $T_w$ 的顶点互不为祖先，所以它们之间没有边（树边也不可能）。

(d)：$u\neq t$，所以 $u$ 不是割点，$G-u$ 连通，$T_v$ 中有一条边伸到 $T_v\cup\{u\}$ 之外。由 (c)，$T_v$ 中的顶点只与它的祖先和后代相邻；它的后代都在 $T_v$ 中，它在 $T_v$ 之外的祖先是 $u$ 与 $u$ 的真祖先。所以这条边连到 $u$ 的一个真祖先；它不是树边（$T_v$ 与外面之间的树边只有 $vu$），是返回边。$u=s$ 时，$s$ 的真祖先只有 $t$。
\end{proof}

\begin{remark}[讲义的笔误]
讲义性质 (b) 写的是“若 $u$ 是 $v$ 的祖先，则 $\texttt{dfn}(u)>\texttt{dfn}(v)$”，不等号反了。(c) 中“$E(V(T_v),E(T_w))=\varnothing$”应为 $E(V(T_v),V(T_w))=\varnothing$。
\end{remark}

\topic{二、low 值}

\begin{definition}\label{def:st-numbering:low}
$v$ 的 \term{low 值}：$\texttt{low}[v]=\min\{\texttt{dfn}[w]: w$ 与 $T_v$ 中某个顶点相邻$\}$，其中“相邻”包括树边 $v\,\texttt{parent}[v]$。\texttt{lowfrom[v]} 记录这个最小值是从哪里得到的：或者是 $v$ 的一个孩子 $c$（$\texttt{low}[v]=\texttt{low}[c]$），或者是 $v$ 的一个邻点 $w$（$\texttt{low}[v]=\texttt{dfn}[w]$）。
\end{definition}

由 (c)，与 $T_v$ 相邻的顶点不是 $T_v$ 中的后代就是 $v$ 的祖先，所以 $\texttt{low}[v]$ 就是“从 $T_v$ 出发用一条边能到达的最高的祖先”的 \texttt{dfn}。由 (d)：$v$ 的双亲 $u\neq s,t$ 时 $\texttt{low}[v]<\texttt{dfn}[u]$；$v=s$ 或 $v$ 是 $s$ 的孩子时 $\texttt{low}[v]=0$。

\begin{remark}
讲义的定义中只算返回边，而讲义的算法（以及下面的程序）把树边 $v\,\texttt{parent}[v]$ 也算了进去，于是总有 $\texttt{low}[v]\leqslant\texttt{dfn}[\texttt{parent}[v]]$。在 $2$-连通图中，由 (d)，树边 $v\,\texttt{parent}[v]$ 只在 $v=s$ 时才可能取到最小值，所以两种定义下 low 值相同（除 $v=s$ 外 \texttt{lowfrom} 也相同），对算法没有影响。
\end{remark}

low 值可以在 DFS 中顺便算出：$v$ 看到一个已访问的邻点 $w$ 时，用 $\texttt{dfn}[w]$ 更新 $\texttt{low}[v]$；$v$ 出栈时，$v$ 的整棵子树都处理完了，用 $\texttt{low}[v]$ 更新双亲的 low 值。（讲义算法 2.9。）

\begin{lstlisting}[language=Python]
def dfn_low(adj, s, t):
    """从边 t-s 开始的 DFS。返回 dfn, parent, low, lowfrom。"""
    n = len(adj)
    dfn, parent = [-1] * n, [-1] * n
    low, lowfrom = [n + 1] * n, [-1] * n
    cur = [0] * n
    dfn[t], dfn[s], parent[s], low[t] = 0, 1, t, 0
    num = 1
    stack = [t, s]
    while stack:
        u = stack[-1]
        if cur[u] == len(adj[u]):          # u 的邻点都看过了：出栈
            stack.pop()
            if u != t:
                v = parent[u]              # 把整棵 T_u 的 low 交给双亲
                if low[v] > low[u]:
                    low[v], lowfrom[v] = low[u], u
        else:
            v = adj[u][cur[u]]
            cur[u] += 1
            if dfn[v] == -1:               # 树边：v 是 u 的孩子
                stack.append(v)
                parent[v] = u
                num += 1
                dfn[v] = num
            elif dfn[v] < low[u]:          # 已访问的邻点：返回边（或通往双亲的树边）
                low[u], lowfrom[u] = dfn[v], v
    return dfn, parent, low, lowfrom
\end{lstlisting}

\begin{example}\label{ex:st-numbering:dfs}
\texttt{Z\_edges=[(0,1),(0,2),(1,2),(1,3),(2,4),(3,4),(3,5),(4,5),(5,0)]}，它是 $2$-连通的。取 $s=1$，$t=0$。邻接表 \texttt{adj=[[1,2,5],[0,2,3],[0,1,4],[1,4,5],[2,3,5],[3,4,0]]}。
\begin{center}
\begin{tikzpicture}[scale=0.95]
  \node[gv] (0) at (0,0) {0}; \node[gv] (1) at (1.5,1.2) {1}; \node[gv] (2) at (1.5,-1.2) {2};
  \node[gv] (3) at (3.2,1.2) {3}; \node[gv] (4) at (3.2,-1.2) {4}; \node[gv] (5) at (4.7,0) {5};
  \foreach \a/\b in {0/1,0/2,1/2,1/3,2/4,3/4,3/5,4/5,5/0} \draw[ge] (\a) -- (\b);
  \node[idx] at (0,-0.55) {$t$}; \node[idx] at (1.5,1.7) {$s$};
\end{tikzpicture}
\end{center}
DFS 的过程（栈从 \texttt{[0,1]} 开始）：
\begin{center}
\small
\begin{tabular}{lll}
\toprule
栈顶 & 看到的邻点 & 结果 \\
\midrule
$1$ & $0$（\texttt{dfn} $0$） & 已访问，$\texttt{low}[1]=0$，$\texttt{lowfrom}[1]=0$ \\
$1$ & $2$ & 新，$\texttt{dfn}[2]=2$，双亲 $1$，栈 \texttt{[0,1,2]} \\
$2$ & $0$ & 已访问，$\texttt{low}[2]=0$，$\texttt{lowfrom}[2]=0$ \\
$2$ & $1$ & 已访问（\texttt{dfn} $1\geqslant\texttt{low}[2]$，不更新） \\
$2$ & $4$ & 新，$\texttt{dfn}[4]=3$，双亲 $2$，栈 \texttt{[0,1,2,4]} \\
$4$ & $2$ & 已访问，$\texttt{low}[4]=2$，$\texttt{lowfrom}[4]=2$ \\
$4$ & $3$ & 新，$\texttt{dfn}[3]=4$，双亲 $4$，栈 \texttt{[0,1,2,4,3]} \\
$3$ & $1$ & 已访问，$\texttt{low}[3]=1$，$\texttt{lowfrom}[3]=1$ \\
$3$ & $4$ & 已访问，不更新 \\
$3$ & $5$ & 新，$\texttt{dfn}[5]=5$，双亲 $3$，栈 \texttt{[0,1,2,4,3,5]} \\
$5$ & $3$ & 已访问，$\texttt{low}[5]=4$，$\texttt{lowfrom}[5]=3$ \\
$5$ & $4$ & 已访问，$\texttt{low}[5]=3$，$\texttt{lowfrom}[5]=4$ \\
$5$ & $0$ & 已访问，$\texttt{low}[5]=0$，$\texttt{lowfrom}[5]=0$ \\
$5$ & （看完） & 出栈；$\texttt{low}[3]=0$，$\texttt{lowfrom}[3]=5$ \\
$3$ & （看完） & 出栈；$\texttt{low}[4]=0$，$\texttt{lowfrom}[4]=3$ \\
$4$ & $5$ & 已访问，不更新；然后看完，出栈（$\texttt{low}[2]=0$ 已经最小） \\
$2$ & （看完） & 出栈 \\
$1$ & $3$ & 已访问，不更新；看完，出栈 \\
$0$ & $1,2,5$ & 都已访问；出栈 \\
\bottomrule
\end{tabular}
\end{center}
结果：
\[
\begin{aligned}
&\texttt{dfn}=[0,1,2,4,3,5],&&\texttt{parent}=[-1,0,1,4,2,3],\\
&\texttt{low}=[0,0,0,0,0,0],&&\texttt{lowfrom}=[-1,0,0,5,3,0].
\end{aligned}
\]
DFS 树是路 $0-1-2-4-3-5$；返回边是 $02$、$13$、$45$、$50$。
\end{example}

\topic{三、st-编号算法}

算法维护一个栈，并把顶点分成\term{老点}（已经进过栈，\texttt{mark}$=1$）与\term{新点}。开始时只有 $s,t$ 是老点，栈为 \texttt{[t,s]}（$s$ 在顶上）。每一步：
\begin{enumerate}
\item 出栈一个顶点 $u$，给它下一个编号。
\item 对 $u$ 的每个是新点的孩子 $v$：从 $v$ 出发沿 \texttt{lowfrom} 往下走，把经过的新点依次标为老点，直到走到一个老点 $w$ 为止。得到路 $u,v,\dots,w$（中间都是刚标的老点）。把中间这些点进栈，使 $v$ 在最上面。
\item 对 $u$ 的每个是新点、并且 $\texttt{dfn}$ 比 $u$ 大的邻点 $v$（即从后代 $v$ 连到 $u$ 的返回边）：从 $v$ 出发沿 \texttt{parent} 往上走，同样标记、进栈，$v$ 在最上面。
\end{enumerate}

\begin{lstlisting}[language=Python]
def st_number(adj, s, t, dfn, parent, lowfrom):
    """2-连通图的 st-编号（讲义算法 2.10）。stn[v] 是 v 的编号 1..n。"""
    n = len(adj)
    stn = [-1] * n
    mark = [0] * n                 # 1 = 老点
    mark[s] = mark[t] = 1
    count = 1
    stack = [t, s]
    while stack:
        u = stack.pop()
        stn[u] = count             # 讲义此处误写为 stn[i]
        count += 1
        for v in adj[u]:           # 第 2 步：树边 u -> v
            if mark[v] == 0 and parent[v] == u:
                path, w = [], v
                while mark[w] == 0:
                    mark[w] = 1
                    path.append(w)
                    w = lowfrom[w]
                while path:
                    stack.append(path.pop())   # 最后一个先进栈，v 在最上面
        for v in adj[u]:           # 第 3 步：返回边 v -> u，v 是 u 的后代
            if mark[v] == 0 and dfn[v] > dfn[u]:
                path, w = [], v
                while mark[w] == 0:
                    mark[w] = 1
                    path.append(w)
                    w = parent[w]
                while path:
                    stack.append(path.pop())
    return stn

def st_numbering(adj, s, t):
    dfn, parent, low, lowfrom = dfn_low(adj, s, t)
    return st_number(adj, s, t, dfn, parent, lowfrom)
\end{lstlisting}

\begin{example}\label{ex:st-numbering:run}
接着\egref[上面的例子]{ex:st-numbering:dfs}。
\begin{center}
\small
\begin{tabular}{p{1.7cm}p{8.6cm}l}
\toprule
出栈，编号 & 第 2、3 步 & 之后的栈 \\
\midrule
$1$，编号 $1$ & 孩子 $2$ 是新点：标 $2$；$\texttt{lowfrom}[2]=0$ 是老点，停。路 $1,2,0$，$2$ 进栈。\newline
返回边：邻点 $3$ 新、$\texttt{dfn}[3]=4>1$：标 $3$；$\texttt{parent}[3]=4$ 新，标 $4$；$\texttt{parent}[4]=2$ 是老点，停。路 $1,3,4,2$，先进 $4$ 再进 $3$。 & \texttt{[0,2,4,3]} \\
$3$，编号 $2$ & 孩子 $5$ 新：标 $5$；$\texttt{lowfrom}[5]=0$ 老，停。路 $3,5,0$，$5$ 进栈。邻点 $1,4$ 都是老点。 & \texttt{[0,2,4,5]} \\
$5$，编号 $3$ & 邻点 $3,4,0$ 都是老点。 & \texttt{[0,2,4]} \\
$4$，编号 $4$ & 邻点都是老点。 & \texttt{[0,2]} \\
$2$，编号 $5$ & 邻点都是老点。 & \texttt{[0]} \\
$0$，编号 $6$ & & \texttt{[]} \\
\bottomrule
\end{tabular}
\end{center}
\texttt{stn=[6,1,5,2,4,3]}，即顺序 $1,3,5,4,2,0$。检查：$3$（编号 $2$）有邻点 $1$（$1$）与 $4$（$4$）；$5$（$3$）有 $3$（$2$）与 $0$（$6$）；$4$（$4$）有 $3$（$2$）与 $2$（$5$）；$2$（$5$）有 $1$（$1$）与 $0$（$6$）。每个中间顶点都有一小一大两个邻点。
\end{example}

\topic{四、为什么算法是对的}

\begin{theorem}\label{thm:st-numbering:correct}
（讲义定理 2.14）设 $G$ 是 $2$-连通图，$st\in E(G)$。则 \texttt{st\_numbering} 在 $O(n+m)$ 时间内给出 st-编号。特别地，每个 $2$-连通图对每条边 $st$ 都有 st-编号。
\end{theorem}

\begin{remark}[证明思路]
关键是两件事。第一，老点的集合始终“对祖先封闭”：一个点是老点，它的祖先都是老点。第二，每次进栈的一条路 $u_1u_2\dots u_k$（$u_1$ 刚出栈，$u_k$ 是原来的老点）中，$u_k$ 还在栈里，没有出栈；于是 $u_2,\dots,u_{k-1}$ 压在 $u_k$ 上面，编号依次落在 $u_1$ 与 $u_k$ 之间，并且依次递增。这样每个 $u_i$ 的两个路上邻点 $u_{i-1},u_{i+1}$ 一个编号小、一个编号大。
\end{remark}

\begin{proof}
\emph{(1) 老点集对祖先封闭。}开始时老点是 $t,s$，成立。第 2 步：$u$ 是老点，$v$ 是它的孩子。沿 \texttt{lowfrom} 走：若 $\texttt{lowfrom}[c]$ 是 $c$ 的孩子，就往下走一层；否则 $\texttt{lowfrom}[c]=w$ 是 $c$ 的一个邻点，$\texttt{dfn}[w]=\texttt{low}[c]=\texttt{low}[v]$（沿孩子走时 low 值不变）。由 (d)，$\texttt{low}[v]<\texttt{dfn}[u]$（$u\neq s$），或 $\texttt{low}[v]=0$（$u=s$），所以 $w$ 是 $u$ 的真祖先，由归纳假设是老点。路上从 $v$ 往下走过的点 $v=c_0,c_1,\dots,c_j$ 是一条树中往下的路；它们原来都是新点（若某个 $c_i$ 是老点，它的祖先 $v$ 也是老点，矛盾），所以路在 $c_j$ 之后才遇到老点 $w$。标完后，新的老点 $c_0,\dots,c_j$ 的祖先是 $u$ 的祖先（老点）与 $c_0,\dots,c_{i-1}$（刚标的），封闭性保持。第 3 步：从后代 $v$ 沿双亲往上走到第一个老点 $x$，把路上的点都标为老点；这些点的祖先是 $x$ 的祖先（老点）与路上的点，封闭性保持。

\emph{(2) 一个顶点出栈并处理完后，它的邻点全是老点。}$u$ 的邻点是它的双亲与祖先（由 (1) 是老点）、孩子（第 2 步之后都是老点）、以返回边相连的后代（第 3 步之后都是老点）。

\emph{(3) 每条进栈的路 $u_1,u_2,\dots,u_k$ 的终点 $u_k$ 不是 $u_1$，并且还在栈中。}第 2 步中 $u_k=w$ 是 $u$ 的真祖先，$\neq u_1$。第 3 步中，往上走时最晚在 $u$ 的孩子处遇到老点（第 2 步已经把 $u$ 的孩子都标成了老点），所以 $u_k\neq u_1$。$u_k$ 是老点，进过栈；若它已经出栈，由 (2) 它的邻点 $u_{k-1}$ 在那时已是老点，而 $u_{k-1}$ 直到这一步才被标记，矛盾。

\emph{(4) 编号。}进栈后，$u_2$ 在最上面，$u_{k-1}$ 在最下面，都在 $u_k$ 之上。栈是后进先出的，以后进栈的点都压在上面，所以 $u_2,u_3,\dots,u_{k-1},u_k$ 按这个顺序出栈，编号递增；$u_1$ 已经编号，比它们都小。于是对 $2\leqslant i\leqslant k-1$，$u_i$ 的邻点 $u_{i-1}$ 编号比它小，$u_{i+1}$ 编号比它大。

\emph{(5) 每个顶点恰好编号一次。}每个顶点至多被标记一次，标记时恰好进栈一次，出栈时编号。若最后还有新点，由 (1) 取一个双亲 $p$ 是老点的新点 $y$；$p$ 进过栈、出过栈，出栈时第 2 步会把孩子 $y$ 标为老点，矛盾。所以所有顶点都进栈、编号，编号是 $1,\dots,n$。$s$ 在栈顶，最先出栈，编号 $1$；$t$ 在栈底，最后出栈，编号 $n$。除 $s,t$ 外每个顶点都是作为某条路的中间点 $u_i$ 被标记的，由 (4) 它有一大一小两个邻点。

\emph{时间。}\texttt{dfn\_low} 是一次 DFS，$O(n+m)$。\texttt{st\_number} 中每个顶点出栈一次，出栈时扫描它的邻接表两遍；沿 \texttt{lowfrom} 或 \texttt{parent} 走的每一步都标记一个新点，总共 $n$ 步。合计 $O(n+m)$。
\end{proof}

讲义的证明给出了 (3)、(4) 的论证，但没有说明 (1) 为什么成立，也没有说明 (2)（“一旦点 $v$ 出栈后，$v$ 的邻点都会变成老点”）为什么成立，而这两点是 (3) 的依据。

\topic{五、low 值的另一个用处：割点}

st-编号要求 $2$-连通。反过来，low 值可以用来找出一个连通图的全部割点，从而判断它是否 $2$-连通。这里用只算返回边的 low 值（不算通往双亲的树边）。

\begin{proposition}\label{prop:st-numbering:cut}
在连通图的 DFS 树（根为 $r$）中：$r$ 是割点当且仅当它有至少两个孩子；$u\neq r$ 是割点当且仅当它有一个孩子 $c$，使 $T_c$ 中没有返回边连到 $u$ 的真祖先，即 $\texttt{low}'[c]\geqslant\texttt{dfn}[u]$，其中 $\texttt{low}'[c]=\min(\{\texttt{dfn}[c]\}\cup\{\texttt{dfn}[w]: T_c$ 中某点与 $w$ 由返回边相连$\})$。
\end{proposition}
\begin{proof}
由 (c)，不同孩子的子树之间没有边，$T_c$ 只通过树边 $cu$ 与返回边连到外面，返回边只连到祖先。$u\neq r$：若某个孩子 $c$ 的 $T_c$ 没有返回边连到 $u$ 的真祖先，删去 $u$ 后 $T_c$ 与 $r$ 不连通；反之，若每个孩子的子树都有返回边连到真祖先，删去 $u$ 后每个 $T_c$ 都连到 $r$ 所在的部分（含 $u$ 的全部真祖先，它们由树边相连），$G-u$ 连通。$r$：孩子的子树之间没有边，有两个孩子时删去 $r$ 不连通；只有一个孩子时 $T-r$ 仍是一棵树。
\end{proof}

\begin{lstlisting}[language=Python]
def cut_vertices(adj, r=0):
    """连通图的全部割点（从 r 出发的一次 DFS）。"""
    n = len(adj)
    dfn, parent, low = [-1] * n, [-1] * n, [0] * n
    children = [0] * n
    cur = [0] * n
    dfn[r] = low[r] = 0
    num, stack, cut = 0, [r], set()
    while stack:
        u = stack[-1]
        if cur[u] < len(adj[u]):
            v = adj[u][cur[u]]
            cur[u] += 1
            if dfn[v] == -1:
                num += 1
                dfn[v] = low[v] = num
                parent[v] = u
                children[u] += 1
                stack.append(v)
            elif v != parent[u]:                 # 返回边（不算通往双亲的树边）
                low[u] = min(low[u], dfn[v])
        else:
            stack.pop()
            p = parent[u]
            if p != -1:
                low[p] = min(low[p], low[u])
                if p != r and low[u] >= dfn[p]:
                    cut.add(p)
    if children[r] >= 2:
        cut.add(r)
    return cut
\end{lstlisting}

（这个程序假定图是简单图：若有重边，“$v\neq$ 双亲”会把与双亲之间的第二条边也排除掉。）

\topic{小结}

\begin{itemize}
\item DFS 树的返回边只连接祖先与后代，$2$-连通性等价于每棵子树都有返回边越过双亲（性质 (d)）。
\item low 值记录子树“最高能连到哪里”，\texttt{lowfrom} 记录这条连接怎么走。
\item st-编号算法每次从刚编号的点出发，沿 \texttt{lowfrom}（往下再跳上去）或 \texttt{parent}（往上）找一条通向老点的路，把路上的点压在那个老点上面。正确性靠两件事：老点集对祖先封闭；路的终点还在栈中。
\end{itemize}

\homework{st-编号的练习}[st-编号]

\exercisepart{练习}

\begin{exercises}{A}{DFS 树与 st-编号}
\item\label{ex:st-numbering:A1} 证明定理~\ref{thm:st-numbering:correct} 的逆命题：设 $n\geqslant3$，$st\in E(G)$。若 $G$ 有 st-编号，则 $G$ 是 $2$-连通的。
\begin{solution}
由定理~\ref{thm:st-numbering:prefix}，$G$ 连通（取 $i=n-1$，前缀与 $\{t\}$ 都连通，$t$ 有编号更小的邻点）。设 $x$ 是割点，$G-x$ 的一个分支 $B$。

若 $s,t\notin B$：取 $B$ 中编号最大的顶点 $b$，它有一个编号更大的邻点，不在 $B$ 中，只能是 $x$；取 $B$ 中编号最小的顶点 $b'$，它有编号更小的邻点，只能是 $x$。于是 $\text{编号}(x)>\text{编号}(b)\geqslant\text{编号}(b')>\text{编号}(x)$，矛盾。

若 $s$ 或 $t$ 在 $B$ 中：$st$ 是边，$x\notin\{s,t\}$ 时 $s,t$ 在 $G-x$ 的同一个分支中；取另一个不含 $s,t$ 的分支，回到上一种情形。$x=s$ 时，取 $G-s$ 中不含 $t$ 的分支 $B$，其中编号最小的顶点要有编号更小的邻点，只能是 $s$；编号最大的顶点要有编号更大的邻点，只能是 $s$，而 $s$ 的编号是 $1$，矛盾。$x=t$ 对称。
\end{solution}

\item\label{ex:st-numbering:A2} 在\egref[例子]{ex:st-numbering:dfs}的图上取 $s=0$，$t=1$，写出 \texttt{dfn}、\texttt{low}、\texttt{lowfrom} 与 st-编号。
\begin{solution}
从 $t=1$ 进入 $s=0$：栈 \texttt{[1,0]}。$0$ 看 $1$（老，\texttt{low[0]=0}，\texttt{lowfrom[0]=1}）、$2$（新，\texttt{dfn}$=2$）；$2$ 看 $0$（\texttt{dfn} $1$，\texttt{low[2]=1}）、$1$（\texttt{dfn} $0$，\texttt{low[2]=0}，\texttt{lowfrom[2]=1}）、$4$（新，$3$）；$4$ 看 $2$（\texttt{low[4]=2}）、$3$（新，$4$）；$3$ 看 $1$（\texttt{low[3]=0}，\texttt{lowfrom[3]=1}）、$4$、$5$（新，$5$）；$5$ 看 $3$（\texttt{low[5]=4}）、$4$（$3$）、$0$（$1$，\texttt{lowfrom[5]=0}）；$5$ 出栈，\texttt{low[3]} 已是 $0$；$3$ 出栈，\texttt{low[4]=0}，\texttt{lowfrom[4]=3}；$4$、$2$、$0$、$1$ 依次出栈。
\[
\begin{aligned}
&\texttt{dfn}=[1,0,2,4,3,5],&&\texttt{parent}=[1,-1,0,4,2,3],\\
&\texttt{low}=[0,0,0,0,0,1],&&\texttt{lowfrom}=[1,-1,1,1,3,0].
\end{aligned}
\]
st-编号：出栈 $0$（编号 $1$）；孩子 $2$ 新，标 $2$，$\texttt{lowfrom}[2]=1$ 老，进栈 $2$；返回边 $5$（$\texttt{dfn}\,5>1$），标 $5$，$\texttt{parent}[5]=3$ 新，标 $3$，$\texttt{parent}[3]=4$ 新，标 $4$，$\texttt{parent}[4]=2$ 老；进栈 $4,3,5$，栈 \texttt{[1,2,4,3,5]}。出栈 $5$（$2$）、$3$（$3$）、$4$（$4$）、$2$（$5$）、$1$（$6$），途中没有新点。\texttt{stn=[1,6,5,3,4,2]}，顺序 $0,5,3,4,2,1$。检查：$5$ 的邻点 $0$（$1$）、$3$（$3$）；$3$ 的邻点 $5$（$2$）、$4$（$4$）；$4$ 的邻点 $3$（$3$）、$2$（$5$）；$2$ 的邻点 $0$（$1$）、$1$（$6$）。
\end{solution}

\item\label{ex:st-numbering:A3} 用 \texttt{cut\_vertices} 找出图 \texttt{[(0,1),(1,2),(2,0),(2,3),(3,4),(4,5),(5,3)]}（两个三角形由一条边 $23$ 相连）的割点，写出 DFS 过程中的 \texttt{dfn} 与 \texttt{low}。
\begin{solution}
邻接表 \texttt{[[1,2],[0,2],[1,0,3],[2,4,5],[3,5],[4,3]]}，从 $0$ 出发：$0\to1\to2$（\texttt{dfn} $0,1,2$）；$2$ 看 $1$（双亲，跳过）、$0$（返回边，\texttt{low[2]=0}）、$3$（新，$3$）；$3\to4$（$4$）$\to5$（$5$）；$5$ 看 $4$（双亲）、$3$（返回边，\texttt{low[5]=3}）。出栈 $5$：\texttt{low[4]=min(4,3)=3}；$p=4\neq0$，$\texttt{low}[5]=3<\texttt{dfn}[4]=4$，$4$ 不是割点。出栈 $4$：\texttt{low[3]=3}；$\texttt{low}[4]=3\geqslant\texttt{dfn}[3]=3$，$3$ 是割点。出栈 $3$：$\texttt{low}[2]=0$；$\texttt{low}[3]=3\geqslant\texttt{dfn}[2]=2$，$2$ 是割点。出栈 $2$：\texttt{low[1]=0}；$0<1$，$1$ 不是割点。出栈 $1$：$p=0$ 是根，不判断。根 $0$ 只有一个孩子。割点为 $\{2,3\}$。最终 $\texttt{dfn}=[0,1,2,3,4,5]$，$\texttt{low}=[0,0,0,3,3,3]$。
\end{solution}
\end{exercises}
